\chapter{Snow basics}
\label{cap:snow_basics}

Snow is a fundamental component of the cryosphere, which encompasses all regions of the Earth where water exists in solid form. 
Alongside sea ice, glaciers, ice caps and ice sheets, snow cover plays an essential role in regulating Earth's climate by reflecting a significant 
fraction of the incoming solar radiation, consequently reducing the amount of energy absorbed by the ground. In addition, snow also acts as a thermal 
insulator due to its low thermal conductivity, protecting the soil and the ecosystem from  the abrupt temperature fluctuations typical of winter. Snow 
is also a key factor in the hydrological cycle: water stored in solid form during winter is gradually released into streamflow through spring and summer, 
ensuring a continuous water supply and mitigating the severity of certain droughts. At the same time, snow has considerable implications for human 
activities. For instance, it plays a vital role in the economy of the Alpine region by supporting winter tourism through skiing and other winter sports. 
Its contribution to runoff is equally important, as it supports irrigation and provides a source of drinking water. Meltwater also provides a 
substantial contribution to hydropower generation. With the ongoing climate change, these dynamics are expected to shift, making the 
understanding and modelling of snow related processes increasingly important. Earlier snowmelt and reduced solid water storage are among the most direct 
consequences of a temperature rise, with potential repercussions for multiple aspects of human civilization. The aim of this chapter is therefore to 
provide a broad overview of snow dynamics, before focusing on snow measurement techniques, which are essential to assess the state of the snowpack.










\section{Snow physics and dynamics}

The formation of snow is the result of a chain of microphysical processes occurring within clouds, from initial ice nucleation to the eventual fall 
of snowflakes to the ground. 

For the atmosphere to release its water content, whether in liquid or solid form, the air must first reach a state of supersaturation with respect 
to water or ice. This condition is typically achieved through upward motion of air masses which causes adiabatic cooling of air parcels. The 
temperature drop strongly reduces air's capacity to hold water vapor, driving the parcel toward, and eventually beyond, saturation \parencite{libro_ATMO}. 

Snowfall can only be produced by clouds that extend above the $\SI{0}{\degreeCelsius}$ level. Although temperatures within these cold clouds can 
drop well below freezing point, supercooled water droplets can still coexist with ice particles. The freezing process, known as ice nucleation, 
occurs through two different mechanisms:
\begin{itemize}[label=$\blacktriangleright$]
    \item \textit{homogeneous nucleation}, which occurs in pure water droplets without foreign particles. The first step in the freezing process 
    is the formation of an ice embryo, which must exceed a critical size to produce a decrease in the free energy of the system. If this condition is met, 
    the freezing process continues until the entire water droplet is turned into ice. The average size of the ice embryo increases as temperature 
    decreases, so that below a certain threshold, which is around $\SI{-40}{\degreeCelsius}$, freezing becomes almost certain. 
    If homogeneous nucleation were the only mechanism responsible for ice particle formation, snowfall would be extremely rare at mid-latitudes. 
    Since this is not the case, an additional mechanism is necessary. 
    \item \textit{heterogeneous nucleation}, which relies on the presence of a foreign particle. Cloud droplets typically contain such impurities 
    acting as freezing nuclei. Their presence promotes the formation of an ice-like structure, thereby enabling freezing at temperatures considerably 
    higher than those required for homogeneous nucleation. It is also possible for foreign particles not to be part of the water droplet itself, 
    but instead to trigger freezing through contact or deposition. Regardless of the specific mechanism involved, the most effective ice nuclei are 
    insoluble in water and have a crystallographic structure similar to ice.  
\end{itemize}

Once nucleated, the ice particle grows through further microphysical processes, whose relative dominance determines its final shape and size. The three 
main crystal growth mechanisms are:
\begin{itemize}[label=$\blacktriangleright$]
    \item \textit{vapor deposition}, which occurs whenever the air is supersaturated with respect to ice. For example, in mixed-phase clouds dominated by 
    supercooled droplets, vapor molecules are deposited onto the ice crystals, causing them to grow significantly faster than the surrounding water 
    droplets. The growth rate reaches a maximum around $\SI{-14}{\degreeCelsius}$, since at that temperature the difference between the saturated 
    vapor pressure over water and over ice is largest, thus driving vapor molecules more effectively towards ice crystals. Studies on the growth 
    of ice crystals by deposition have shown a strong relationship between crystal shape and the temperature at which the growth process takes place. 
    This strong temperature dependence produces an alternating pattern between plate-like and column-like ice particles, as shown in Figure \ref{fig:snow_flakes}.
    \item \textit{riming}, which occurs exclusively in mixed-phase clouds. Whenever an ice crystal collides with a supercooled droplet, the latter may 
    freeze onto the former, resulting in an increase in ice mass. The collision frequency depends on cloud density and composition, but typically 
    numerous riming events take place, sometimes producing a graupel, namely an ice particle in which the original ice crystal can no longer be identified. 
    Riming can also lead to the formation of hailstones, which typically range from a few millimeters to several centimetres in diameter.
    \item \textit{aggregation}, caused by collisions between ice particles. The probability of adhesion is strongly linked to crystal shapes and increases 
    with increasing temperature, reaching a maximum at around $\SI{-5}{\degreeCelsius}$, where ice surfaces become particularly "sticky".
\end{itemize}

\begin{figure}[htbp]
    \centering
    \includegraphics[width=0.8\linewidth]{snow_basics/snow_flakes.png}
    \caption{The Nakaya diagram illustrates the morphology of ice crystals growing via vapor deposition as a function of cloud's thermodynamic state. 
    Crystal shape is primarily governed by temperature, which determines whether it develops in a planar or columnar structure, while the degree of 
    supersaturation with respect to ice modulates the complexity of the resulting form, with higher supersaturations generally linked to more 
    complex structures. Image taken from \cite{snow_crystals}.}
    \label{fig:snow_flakes}
\end{figure}

Once an ice particle grows large enough to overcome updraft forces and turbulent motion within the cloud, it begins to fall toward the ground. The 
temperature profile encountered by the descending particle ultimately determines the type of precipitation that reaches the surface. For 
snowfall to occur, the snowflake, composed of up to hundreds of ice crystals, must remain predominantly below freezing and avoid melting completely
before reaching the ground. Typically, snowflakes have fall rates of approximately $\SI{1.5}{\meter\per\second}$. Snowfall amounts are measured either 
by depth or by snow water equivalent (SWE), defined as the height of the water column resulting from a complete melting of the snow. Snow depth typically 
increases of approximately $\SI{1}{\centi\meter\per\hour}$ during snowfall events. Newly fallen snow density ranges from $\SI{20}{\kilo\gram\per\meter\cubed}$ 
to $\SI{300}{\kilo\gram\per\meter\cubed}$ depending on wind patterns and weather conditions, with higher temperatures generally leading to higher 
density values \parencite{libro_SNOW}.



\subsection{Snow metamorphism and compaction}
Accumulation of snowflakes on bare ground or atop older snow results in the formation of a new snow layer. Snow crystals form a solid matrix 
delimiting pores filled with humid air or liquid water or both, depending on layer temperature. This structure can be clearly appreciated in Figure 
\ref{fig:snow_porous}, where the coexistence of the solid, liquid and gaseous phases of water is visible. Snow is thus a porous medium, whose 
properties are closely linked to its texture.

The arrangement of snow crystals evolves over time as a result of thermodynamic interactions among the constituents of the snow layer. This process, 
known as \textit{snow metamorphism}, is strongly influenced by thermal and meteorological conditions. Since each layer within the snowpack is subject 
to different burial depths, temperature gradients and moisture presence, it follows a distinct metamorphic pathway. Consequently, the snowpack develops 
a stratified structure, with individual layers exhibiting different crystal characteristics.

A key factor governing the evolution of ice particles is the presence or absence of liquid water, which motivates the distinction between dry and wet 
snow metamorphism. The former occurs when the pores are filled only by air, typically saturated with water vapor. The main mechanisms underlying dry 
snow metamorphism are:
\begin{itemize}[label=$\blacktriangleright$]
    \item \textit{equilibrium growth}, driven by crystal shape. Local variations in surface curvature generate corresponding variations in 
    equilibrium vapor pressure, inducing local pressure gradients that drive vapor from convex toward concave (or less convex) regions. Over time this 
    transfer, sustained by sublimation on convex areas and deposition on concave ones, rounds ice grains, driving snow crystals toward a spherical shape.
    \item \textit{kinetic growth}, related to the presence of a temperature gradient. Because saturation vapor pressure depends strongly on 
    temperature, such a gradient causes vapor to diffuse from the warmest crystals toward the coldest ones. Here too, sublimation on warm crystals and 
    deposition on cold ones sustain the process, which, when the temperature gradient is strong enough, leads to the formation on facets.
\end{itemize}

\begin{figure}[htbp]
    \centering
    \includegraphics[width=0.7\linewidth]{snow_basics/snow_porous.png}
    \caption{Thin section of wet snow under transmitted light, where ice grains delimit pore spaces containing both liquid water and air. The presence 
    of air-filled pores suggest that the sample can potentially retain additional water. Image from \cite{snow_slice}}
    \label{fig:snow_porous}
\end{figure}

These two mechanisms are rarely mutually exclusive: they act simultaneously within the snowpack, with their relative strengths determined 
mainly by the magnitude of the local temperature gradient compared to the grain's curvature effect. In particular, the kinetic growth of snow 
particles is a matter of considerable concern in avalanche forecasting, since faceted crystals, if the temperature gradient persists over time, can 
evolve in depth hoar. This type of layer is notoriously weak and prone to failure, often resulting in large-scale avalanches.

When liquid water is present within the snowpack, the dynamics change substantially: mass exchanges between the three water phases become essential to 
understand wet snow metamorphism. The main driver of ice particles evolution is, once again, crystal surface curvature, with convex areas associated with 
lower melting point temperatures. This topological property leads to the melting of the most convex surfaces and to refreezing on the more concave ones, 
driving ice particles towards the equilibrium spherical shape. The efficiency of this metamorphism channel is linked to the freedom water has to move 
within the snowpack: at very low water content mass exchanges between the solid and gaseous phases may still dominate.

Regardless of the specific metamorphic process involved, the evolution of ice crystals ultimately results in the compaction of the snow layer. Wind also 
plays a crucial role in increasing snow density, primarily through the fragmentation of snowflakes upon collision, which allows for a tighter packing of 
ice particles. Finally, gravity contributes to snow compaction as well, since buried snow layers must withstand an increasing load exerted by the weight 
of the overlying layers, resulting in higher deformation. The combined effects of metamorphism and deformation strain are typically described 
through a Newtonian viscosity approach. Snow viscosity depends on several parameters, including temperature, liquid water content and porosity and its
values are generally retrieved by means of empirical relations. Overall, the density of older snow layers can easily exceed values of $\SI{400}{\kilo\gram\per\meter\cubed}$, 
sometimes even reaching $\SI{500}{\kilo\gram\per\meter\cubed}$ in specific locations, depending on water content and time of the year \parencite{rho_max}.



\subsection{Snowpack energy balance}

The evolution of a snowpack throughout the season is governed not only by the microphysical processes described above, but also by the exchange of 
energy between the snow cover, the atmosphere and the underlying ground. Neglecting horizontal energy transfers, the internal energy of the snowpack 
evolves in time as:
\begin{equation}
    \frac{d\mathcal{H}}{dt}\,=\,R_N\,+\,H_S\,+\,H_L\,+\,H_P\,+\,G,
    \label{eq:energy_balance}
\end{equation}
where $\mathcal{H}$ is the internal energy per unit area and $R_N$ is the net radiative flux at the surface. $H_S$ and $H_L$ denote, respectively, 
the turbulent fluxes of sensible and latent heat, governed by the action of turbulent eddies within the boundary layer. $H_P$ accounts for the 
energy exchange associated with precipitation, since rainfall carries additional heat that must be distributed within the snowpack. Finally, $G$ 
represents the energy exchange with the ground, which is especially important during the deposition of the first snow layer. Each term in Equation 
\eqref{eq:energy_balance} is expressed in $\SI{}{\watt\per\meter\squared}$ and, by convention, a positive sign indicates an energy flux directed into 
the snow layer. 

To fully understand the term $R_N$, it is essential to consider the dual nature of radiation in the atmosphere. Two distinct sources of radiation can in 
fact be identified: the Sun, with a surface temperature of around $\SI{5800}{\kelvin}$, and the Earth-atmosphere system, whose temperature 
ranges from $\SI{200}{\kelvin}$ up to $\SI{350}{\kelvin}$. According to the Wien's law, the two blackbody spectra peak respectively at roughly 
$\SI{0.5}{\micro\meter}$ and $\SI{10}{\micro\meter}$ with negligible overlap, thereby justifying a decoupled treatment of the two spectral 
domains. Solar radiation is commonly referred to as \textit{shortwave} (SW), while the latter is known as \textit{longwave} (LW). Therefore, the term 
$R_N$ can be explicitly written as
\begin{equation}
    R_N\,=\,SW\downarrow\,-\,SW\uparrow\,+\,LW\downarrow\,-\,LW\uparrow,
    \label{eq:net_radiation}
\end{equation}
where a distinction is made between incoming and outgoing fluxes, as denoted by the direction of the arrows. 

Incoming shortwave radiation is subject to absorption and scattering by atmospheric constituents. Consequently, the total SW radiation that reaches 
the Earth's surface is the sum of a direct beam transmitted through the atmosphere and a diffuse component, resulting from multiple reflection and 
refraction events. The relative strengths of these two components is strongly dependent on meteorological and seasonal conditions. For example, at 
the lower solar altitudes typical of winter, the diffuse fraction is larger, since the solar beam must travel through a thicker atmospheric layer, 
increasing scattering probability. The same occurs at dawn or sunset, when the Sun is close to the horizon. Cloud cover also plays a significant 
role: under overcast conditions, up to $80 \%$ of $SW\downarrow$ can occur as diffuse radiation \parencite{snow_hydrology}. In the shortwave spectral 
range, snow features an exceptionally high reflectance, thanks to the combination of its microstructure with the optical properties of ice. The ratio 
between $SW\uparrow$ and $SW\downarrow$, commonly referred to as albedo $\alpha$, typically exceeds $0.9$ for freshly fallen snow, and decreases as 
impurities are deposited on the snow surface. This parameter is rather important, as it controls the amount of shortwave radiation absorbed, which 
generally represents the main source of energy for snowmelt. By albedo definition, this quantity can be calculated as
\begin{equation}
    SW_{net}\,=\,SW\downarrow \left(1\,-\,\alpha\right).
    \label{eq:SW_albedo}
\end{equation}
Even at advanced stages of the melt season, snow albedo remains considerably higher than that of the snow-free surroundings. As the snowpack 
retreats, this bright surface is replaced by darker ground, causing a significant $\alpha$ drop. This results in increased SW absorption, which 
causes further warming and, in turn, accelerates snowmelt itself in a self-reinforcing loop known as snow-albedo feedback.

Snow behaviour in the thermal infrared range, however, differs substantially from that in the shortwave domain. The snowpack absorbs nearly all incoming longwave 
radiation, emitted by atmospheric gases, clouds and the surrounding terrain, while radiating according to the Stefan-Boltzmann law and thus behaving as 
an almost perfect blackbody, with emissivities ranging from $0.97$ to $0.99$. Longwave radiation exchanges occur both during the day and at 
night, and are strongly influenced by meteorological conditions. In fact, water vapor significantly absorbs infrared light and, by  
Kirchhoff's law, is therefore also a strong emitter. As a result, clouds typically increase the downward longwave flux reaching the snow surface, thereby 
reducing the snowpack's radiative cooling capacity. This effect is especially pronounced at night, when the absence of the shortwave component would 
otherwise favor significant longwave losses.

Thanks to its relatively low thermal conductivity, the snowpack also shields the ground from rapid atmospheric temperature changes. This property is 
closely tied to snow microstructure and therefore evolves with metamorphism, as higher snow densities are associated with increased conductivity. 



\subsection{Snow melt and water runoff}

The variations in internal energy $\mathcal{H}$, driven by the heat fluxes listed above, ultimately modify the thermal state of the 
snowpack and determine whether mass is lost through melting or sublimation. The snowpack temperature profile varies on both seasonal and daily scales, 
typically showing large gradients near the snow surface as a consequence of the low thermal conductivity of snow. The temperature oscillations of 
the first layers are damped rather quickly with depth, as the thermal state of the deeper bulk remains relatively stable. During winter, snowpacks 
typically exhibit subfreezing conditions, which poses a significant limitation to runoff generation. In fact, even when the surface temperature rises 
to $\SI{0}{\degreeCelsius}$ and the energy balance is positive, the water produced by surface melting refreezes in deeper layers after some percolation. 
Although this process does not involve any mass loss, it effectively transports latent heat, warming the snowpack and reducing its cold content, which is 
the energy required to bring the whole snowpack to $\SI{0}{\degreeCelsius}$. This isothermal profile is eventually reached and, if a positive energy 
balance still persists, further melting contributes to filling the pores between ice crystals. This liquid water holding capacity is commonly reported 
ranging between $3\%$ and $5\%$ of the snowpack mass, although the exact value depends on the snow's porosity and microstructure. 
Once the water holding capacity is exceeded, the excess water percolates through the snowpack in a process that is 
mainly gravity driven. Because a typical snowpack is stratified into layers of contrasting grain size and permeability, water flow is rarely uniform 
as it tends to concentrate along preferential pathways known as flow finger, which are usually triggered by inhomogeneities in microstructure. 
Typical percolation velocities span more than an order of magnitude, ranging from $\SI{1}{\centi\meter\per\hour}$ 
up to $\SI{20}{\centi\meter\per\hour}$. Ground conditions represent another critical factor in runoff evaluations, since meltwater can in principle 
infiltrate the soil. If the underlying surface is impermeable, such as frozen ground or bedrock, the excess water instead accumulates and flows downward 
along the base of the snow cover. If the snowpack is not isothermal throughout, part of this water may still refreeze, further warming the surroundings. 
The interplay between thermal conditioning, liquid water retention, and preferential flow pathways ultimately governs the timing, volume and rate of 
snowmelt runoff entering local stream networks. 

In high elevation alpine catchments, or at high latitudes, some snow may survive melting through summer. If accumulation persists over multiple years, 
continued densification ultimately leads to the formation of firn, an intermediate stage between snow and ice characterized by densities ranging 
from $\SI{400}{\kilo\gram\per\meter\cubed}$ to $\SI{830}{\kilo\gram\per\meter\cubed}$ and by a porous structure that still permits water percolation. 
As compaction proceeds further, most pores become sealed, marking the transition to glacier ice proper. The rate of this transformation depends on 
both accumulation amounts and temperature, since colder climates results in slower crystal evolution. This process directly links seasonal snow cover 
to the longer lived components of the cryosphere, such as glaciers and ice sheets \parencite{snow_to_ice}. 










\section{Snow observation and measurement}
A detailed knowledge of the snow cover and of its physical properties is a prerequisite for studying snowpack dynamics and for issuing reliable avalanche 
forecasts. Among the quantities routinely collected during snow measurement campaigns, snow depth (HS) plays a central role: it is defined as the vertical 
distance between the top of the snowpack and a reference surface, which for seasonal snowpacks is normally the ground. HS should not be confused with snow 
thickness (DS), which is conventionally the perpendicular distance to the same reference surface. Provided that the slope angle $\phi$ of the measurement 
site is known, the two quantities are related by 
\begin{equation}
    \text{DS}\,=\,\text{HS} \cos{\left(\phi\right)}, 
    \label{eq:hs_to_ds}
\end{equation}
so that the two definitions coincide on flat terrain ($\phi\,=\,0$) and progressively diverge as the slope steepens. The geometrical meaning of Equation 
\eqref{eq:hs_to_ds} is illustrated in Figure \ref{fig:depth_vs_thicness}. Both quantities are expressed in $\SI{}{\centi\meter}$ and are usually rounded 
to the 
\begin{figure}[H]
    \centering
    \includegraphics[width=0.7\linewidth]{snow_basics/depth_vs_thicness.png}
    \caption{Relationship between snow depth (HS) and snow thickness (DS) on a slope. Image adapted from \cite{snow_instrument}.}
    \label{fig:depth_vs_thicness}
\end{figure}
nearest integer.
While these metrics describe only the geometrical extent of the snowpack, other quantities provide deeper insight into its physical state. The most 
hydrologically relevant of these is snow water equivalent (SWE), namely the amount of water stored within the snowpack. Being closely related to snow 
density, it can be expressed as 
\begin{equation}
    \text{SWE}\,=\,\int_{0}^{\text{HS}}\rho_{s}\left(z\right) dz, 
    \label{eq:swe_density_link}
\end{equation}
where $\rho_s\left(z\right)$ is the vertical density profile of the snowpack. The result is given in $\SI{}{\kilo\gram\per\squared\meter}$, and can be 
converted in $\SI{}{\meter}$ (or $\SI{}{\milli\meter}$) by dividing by the water density $\rho_w$. Although formally correct, Equation \eqref{eq:swe_density_link} 
is of limited practical use, since continuous knowledge of the density profile is generally unattainable. What can realistically be measured is the 
density of the individual layers, so that SWE is more commonly computed as 
\begin{equation}
    \text{SWE}\,=\,\sum_{i=1}^{N}\rho_{i}\left(z\right) \Delta z_i, 
    \label{eq:swe_density_layers}
\end{equation}
where $N$ is the number of layers and $\Delta z_i$ is the vertical extent of the $i$-th layer. 

Historically, snow metrics were recorded manually, requiring an operator to carry out manual measurements on a daily basis whenever feasible. 
To ensure continuous data collection even in remote areas, automated weather stations (AWS) have been introduced and are progressively 
replacing manual observations as well. An automated monitoring network provides a comprehensive overview of snow conditions, but achieving high spatial 
resolution remains challenging due to high installation costs and maintenance issues. Consequently, certain data can only be obtained through remote 
sensing, by means of Earth-observing satellites.

While remote sensing provides spatially distributed data across entire regions, manual and automated measurements are taken at specific locations, 
which must be selected according to strict criteria. To ensure the reliability of these in situ observations, the monitoring site must match the 
features of the surrounding terrain. In alpine regions, monitoring sites must be free of large boulders and protected from excessive sun 
exposure, which can cause premature snowmelt and thus yield unrepresentative data. It is equally important to avoid locations prone to preferential 
snow accumulation driven by local wind patterns \parencite{snow_instrument}.



\subsection{Manual snow measurements}

Manual snow depth (HS) records are an effective tool for assessing long-term trends, as they can span several decades. Their longevity is 
due to a straightforward measuring procedure, requiring only a stake or an extendible graduated rod. The snow stake should be vertical, painted 
white to minimize melt and fixed to the ground before accumulation begins. Snow depth observations should be made from a sufficient distance to 
avoid any deformation of the snowpack. Because the stake acts as an obstacle to wind and snow redistribution, localized mounds can form around it. 
Under such circumstances, operator must visually estimate the undisturbed snow surface to recover the actual HS value. For a non-fixed observation 
carried out with a graduated rod, the snow depth should be evaluated at the same location each time. The observer needs to ensure that the device, 
manually inserted into the snowpack, reaches the base surface without penetrating the soil, which would result in an overestimate of the actual HS. 
The rod must be vertical, and an estimate of the slope angle should be paired with every measurement.

Closely linked to snow depth is snowfall depth (HN), which can only be obtained manually. To avoid wind-induced bias this measurement should 
be carried out in a sheltered location, so that snow accumulates undisturbed on an artificial surface for the prescribed interval (usually 24 hours). 
Just before clearing the surface, thereby starting the next observation period, a ruler is inserted vertically into the freshly formed layer to 
obtain a depth measurement.

Although HS and HN define the vertical extension of the snowpack, hydrological applications rely mainly on the snow water equivalent (SWE): the 
two most common methods for manual measurement are snow pits and snow tubes. A snow pit consists of a hole manually dug into the snowpack down to a 
reference surface, usually the ground. Initially, the observer performs a stratigraphic analysis to gain deeper insight into the internal structure 
of the snowpack. The actual SWE measurement is conducted by inserting a sampler of known volume into the exposed snow face; the sample 
is then weighed to assess snow density. This procedure is usually repeated at multiple depths, either for each identified layer or at fixed spatial 
intervals, such as every 20 centimetres. This method enables the observer to reconstruct the snow density profile which, when combined with depth 
information, yields both the layer-by-layer and total snow water equivalent at that specific location. Clearly, the snow pit technique is disruptive 
and needs to be performed at a different site each time. 

If the measurement campaign is focused on a spatial survey rather than a detailed, locally 
defined inspection, the logical choice is to use a snow tube, as it allows for rapid sampling of the entire snow column to directly retrieve the total 
SWE, sacrificing layer-specific information. The snow tube must be inserted vertically into the snowpack until it reaches the ground. An estimate of 
the extracted core volume is made, and the average density is then computed by weighing the snow column. The conversion to SWE is then straightforward. 
This device should be used where the snow water equivalent is less spatially homogeneous, such at locations with heterogeneous exposure to both wind 
patterns and solar radiation.



\subsection{Automatic snow monitoring}

Networks of automatic weather station (AWS) require no manual intervention beyond routine maintenance and instrument calibration, and can provide 
continuous observations even in remote locations where human operation would be impractical. An AWS is typically equipped with a wide range 
of sensors, such as:
\begin{itemize}[label=$\blacktriangleright$]
    \item \textit{anemometer}, which measures wind speed. The blowing wind causes a small windmill, or a structure consisting of three cups mounted 
    on a spindle, to rotate at a rate proportional to wind intensity. This quantity is particularly important, as empirical relations link it to 
    turbulent heat fluxes, thus playing a significant role in energy balance considerations.
    \item \textit{wind wane}, typically coupled with the anemometer to assess wind direction. It consists of a horizontal arm bearing a flat vertical 
    vane at one end. Incoming wind gusts push against the vain, causing the arm to rotate freely until it aligns with the wind direction.
    \item \textit{pyranometer}, which measures shortwave solar radiation. An AWS is typically equipped with both and upward and a downward facing 
    pyranometer, in order to quantity incoming and outgoing SW radiation, thus enabling albedo estimates.
    \item \textit{pyrgeometer}, which measures the longwave (infrared) radiation. Pyrgeometers likewise come in oppositely facing pairs, so that, 
    combined with a pair of pyranometers, an evaluation of the net radiation balance of the surface is possible.
    \item \textit{barometer}, used to evaluate the atmospheric pressure, a fundamental quantity for characterizing the dynamical state of the atmosphere.
    \item \textit{thermo-hygrometer}, designed to simultaneously measure both air temperature and relative humidity at a given height above 
    the ground (typically two meters).
    \item \textit{rain gauge}, consisting of a circular collector and a funnel that channels the collected precipitation into a measuring device. In 
    regions where solid precipitation is possible, the rain gauge may be heated in order to also evaluate the water equivalent of snowfall.
\end{itemize}

In alpine catchments and at high latitudes, these sensors are typically complemented by additional instruments specifically dedicated to snow monitoring.
Mounted together on the same tripod or mounting structure, as shown in Figure \ref{fig:station_close_up}, they constitute what is referred to as an 
alpine automatic weather station, a key tool for the study of snow related dynamics.

Automated HS measurements can be performed using instruments based on either sonic or optical (laser) technology. Sonic sensors transmit an 
ultrasonic pulse towards the snow surface and detect the returning echo. Since the speed of sound depends on meteorological conditions, a concurrent 
temperature reading is required to ensure accuracy. Laser instruments, by contrast, determine the distance to the target surface through 
interferometric analysis of the returned signal. This approach yields significantly higher precision, but at a grater energy cost. 

In both cases, the sensor measures the distance to the upper surface of the snowpack, from which the snow depth is then derived as
\begin{equation}
    \mathrm{HS} = \mathrm{Zero\ depth\ distance} - \mathrm{Distance\ to\ target}.
    \label{eq:auto_HS}
\end{equation}
The zero depth distance must be determined before snow accumulates and constitutes a key step in instrument calibration. To ensure reliable 
measurements, the sensor should be positioned well above the maximum expected snow depth, preventing burial and the consequent loss of data during 
winter. Minimizing interference from 
\begin{figure}[H]
    \centering
    \includegraphics[width=0.9\linewidth]{snow_basics/aws.png}
    \caption{The AWS UPE\_L which is part of an automated network of the Programme for Monitoring of the Greenland Ice Sheet 
    (PROMICE). The constituents are respectively: (1) radiometer LW + SW, (2) inclinometer, (3) satellite antenna, (4) anemometer and wind vane, 
    (5) snow depth sensor, (6) thermometer and hygrometer, (7) pressure transducer, (8) solar panel, (9) data logger and GPS, (10) battery box and 
    (11) ground temperature sensor. Image adapted from \cite{image_aws}.}
    \label{fig:station_close_up}
\end{figure}
the support structure is equally important, since the mast itself acts as a wind obstacle and may induce 
measurement artefacts. For this reason, HS sensors are typically mounted on a horizontal arm stretching outward from the pole, keeping the target 
area sufficiently far from the mounting structure.

SWE monitoring can likewise be automated, with the most common technique relying on snow pillows. This device consists of a rubber bladder filled with 
an antifreeze liquid mixture and connected to a pressure transducer, which yields the SWE estimate. The pillow is installed flush with the ground 
surface, so that it supports the full weight of the overlying snowpack. As snow accumulates, the resulting increase in hydrostatic pressure is 
transduced into snow water equivalent values. Snow pillows were first introduced in the late 1960s, and have proved to be a valuable 
tool for SWE measurements \parencite{SnowPillow_Hood}. However, they are now being progressively replaced by snow scales due to the toxicity of the 
antifreeze fluid, which poses an environmental hazard to the installation site.



\subsection{Remote sensing}

Satellite remote sensing (SRS) serves as a useful tool to fill critical gaps in current knowledge regarding climatic variables, enabling researchers 
to monitor, usually on a daily basis, even remote locations that would otherwise be inaccessible to a human operator. Furthermore, SRS data substantially 
improve meteorological reanalysis products, which constitute a cornerstone of contemporary climate change research. Earth observing satellites can be 
classified into two distinct categories based on their orbital characteristics:

\begin{itemize}[label=$\blacktriangleright$]
    \item Polar orbit satellites, which travel roughly from north to south over the poles, targeting a specific longitudinal strip per pass. The 
    entire globe is thus mapped over consecutive orbits. These satellites typically operate at an altitude of approximately 700 km, completing a full 
    revolution in under two hours.
    \item Geostationary satellites, which occupy a geosynchronous orbit, thereby allowing for real-time meteorological monitoring on a specific target 
    area. Their altitude of approximately 36000 km above Earth's equator results in a significantly lower spatial resolution than that of 
    polar orbit satellites.
\end{itemize}

The choice of orbit is strictly driven by observational priorities. Consequently, satellites dedicated to cryospheric research favor polar orbit, as
maximizing spatial resolution is essential to capture the highly variable distribution of snow coverage across complex terrains. 

Satellite-based observations of snow cover extent (SCE) began in 1967, through the efforts of National Oceanic and Atmospheric Administration (NOAA). 
This monitoring program focused on the Northern Hemisphere, where most of the world's terrestrial snow cover is found, with observations conducted 
on a weekly basis. Despite the preliminary nature of this work, it highlighted the strong latitudinal dependence of snow cover duration and paved 
the way for more comprehensive analyses. One of the main limitations of these early observations was their coarse spatial resolution of 
$\sim \SI{200}{\km}$, which was insufficient to adequately characterize snow cover in many mountainous regions \parencite{NOAA_first_SRS}. 
A significant advance in sensor technology came with the introduction of the Advanced Very High Resolution Radiometer (AVHRR), which measures Earth 
radiance in five spectral bands to derive information on temperature, albedo and cloud coverage distribution. The first tests of a spaceborne sensor 
took place in 1978, but reliable snow cover data have only been available since 1981. This product has been used to estimate snow cover duration 
and trends since the mid-eighties in various mountainous regions such as the Alps, where SRS data reveal significantly negative trends throughout 
the 1990s, followed by a partial recovery at the beginning of the new millennium \parencite{AVHRR_example}. Although AVHRR offers a substantial  
improvement in spatial resolution, with pixels of $\sim \SI{1}{km}$, its data reliability is limited by the presence of only one thermal channel 
for snow-cloud discrimination \parencite{MODIS_GAR}.
To overcome these limitations, the Moderate Resolution Imaging Spectroradiometer (MODIS) was launched in 1999 aboard the Terra satellite, followed by 
a second unit aboard Aqua in 2002. Both instruments operate in polar orbit, jointly covering the entire globe every one to two days. MODIS is a 
36-band spectral sensor spanning from the visible ($\SI{0.41}{\micro m}$) to the long-wave infrared ($\SI{14.4}{\micro m}$), with a spatial resolution 
of $\SI{500}{m}$. To detect the presence of snow on the ground, the Normalized Difference Snow Index (NDSI) is used, defined as
\begin{equation}
    \mathrm{NDSI}\,=\,\frac{\mathrm{G}\,-\,\mathrm{SWIR}}{\mathrm{G}\,+\,\mathrm{SWIR}},
    \label{eq:NDSI}
\end{equation}
where G and SWIR denote the green and shortwave infrared spectral bands respectively. This normalized difference approaches one for a snow-covered 
pixel, since snow exhibits very high reflectance in the visible portion of the electromagnetic spectrum, while strongly absorbing radiation in the 
SWIR. An NDSI value greater than zero indicates the presence of snow in the pixel; in practice, a threshold of 0.4, corresponding to roughly 50 \% 
snow coverage, is commonly adopted to classify a pixel as snow covered \parencite{NDSI_threshold}.


Although MODIS has proven to be a robust and well-validated tool for characterizing snow cover, the information it provides remains partial, 
as it is limited to detecting the presence or absence of snow on the ground. In order to assess how much water is stored in solid form, three main 
spaceborne techniques have been developed: 
\begin{itemize}[label=$\blacktriangleright$]
    \item \textit{passive microwave radiometry (PMW)}, which is by far the most common method to remotely evaluate SWE. Since the snowpack attenuates 
    the microwaves emitted by the Earth, comparing signal intensities between two frequencies ($\sim \SI{18}{\giga\hertz}$, and $\sim \SI{37}{\giga\hertz}$) 
    makes it possible to estimate snow water equivalent. This is the only spaceborne technique that allows for long term studies, as global daily 
    maps have been available since the launch, in 1978, of NASA's Scanning Multichannel Microwave Radiometer (SMMR) \parencite{remote_SWE}. Altough  
    the temporal coverage of these series is considerably better than that of other available sources, several factors affect data reliability, 
    such as signal saturation caused by a deep snowpack.
    \item \textit{gravimetric methods}, which offer a different approach focusing on mass changes. The accumulation of snow on the Earth's surface 
    causes variations in its gravitational field, which can be converted into SWE estimates. This approach is mainly employed by the Gravity Recovery 
    and Climate Experiment (GRACE) mission. Owing to its global coverage, this SWE dataset features a coarse spatial resolution of roughly 
    $\SI{300}{\kilo\meter}$, which nonetheless remains useful for continental scale hydrological studies \parencite{remote_SWE}.
    \item \textit{Synthetic Aperture Radar (SAR)}, which consists in the transmission of electromagnetic waves towards Earth and the analyses of the 
    backscattered echoes from the Earth's surface. SAR can penetrate cloud cover and acquire data under all weather conditions. Snowpack properties, 
    such as SWE and snow grain size can be estimated remotely, as longer SAR wavelengths are even capable to penetrate into the snowpack. Furthermore, 
    the phase information recorded by SAR can be exploited to monitor ground deformation and assess the stability of terrain features. Although SAR 
    provides a spatial resolution on the order of tens of metres, its temporal resolution is significantly lower than that of PMW and gravimetric methods, 
    as the revisit time of a specific location is typically five days \parencite{remote_SAR}.

\end{itemize}









\section{Climate change}

% Citare UNFCCC
Climate change is conventionally defined as an alteration of the climate attributed, directly or indirectly, to human activities that modify the 
composition of the global atmosphere. This anthropogenic signal is superimposed on the natural variability of the climate system and must therefore 
be disentangled from a background of spontaneous fluctuations. Over the past few million years, such fluctuations have produced a cyclic alternation 
between glacial periods, lasting roughly \num{100000} years, and warmer interglacial conditions. These climatic oscillations are caused by periodic 
variations in the orbital parameters governing Earth's revolution around the sun, which generate a weak radiative forcing that is subsequently 
amplified by snow albedo and greenhouse gas feedbacks \parencite{iceage_orbit}. The last glacial period ended around \num{11700} years ago, marking 
the onset of the Holocene, the geological epoch we live in today.

The ongoing global warming, however, cannot be attributed to natural variability, as it is primarily driven by the large amounts of greenhouse gases 
released into the atmosphere by human activities. These gaseous compounds, of which carbon dioxide ($\text{CO}_2$) is the best known, partially absorb 
the longwave radiation emitted by Earth's surface and re-emit it in all directions, so that only a fraction of this energy escapes, thus causing an 
increase in surface air temperature. Since the Industrial Revolution the atmospheric $\text{CO}_2$ concentration has risen steadily, from about 
\num{278} part per million (ppm) in pre-industrial times to $\SI{422.8}{ppm}$ in 2024 \parencite{co2_conce}. During the last 60 years, $\text{CO}_2$ 
concentrations have grown nearly 100 times faster than during previous natural increases, such as the one that occurred at the end of the last ice 
age. As visible in Figure \ref{fig:co2_and_emissions}, this rise is closely related to the total carbon emissions caused by human activities. 
\begin{figure}[H]
    \centering
    \includegraphics[width=0.9\linewidth]{snow_basics/co2_and_emissions.png}
    \caption{Visual superposition of $\text{CO}_2$ concentration in the atmosphere and cumulative $\text{CO}_2$ emissions made by mankind. 
    Image adapted from \cite{co2_conce}.}
    \label{fig:co2_and_emissions}
\end{figure}

Relative to the 1850--1900 climatology, the global mean surface air temperature observed over the 2010--2019 period was about $\SI{1.1}{\degreeCelsius}$ 
higher, and anthropogenic forcing alone is estimated to account for $\SIrange{0.9}{1.3}{\degreeCelsius}$ of this increase, with the best estimate 
being $\SI{1.07}{\degreeCelsius}$ \parencite{atro_forcing}, thus explaining the vast majority of the measured signal. The warming has also accelerated 
significantly in recent years, as the mean increase rate of roughly $\SI{0.18}{\degreeCelsius}$ per decade shown during the 1970--2010 period has 
nearly doubled between 2010 and 2023 \parencite{temp_trend}. This behaviour is clearly visible in Figure \ref{fig:temperature_trend}, which shows 
the time evolution of the global surface temperature anomaly with respect to the 1880--1920 mean.
\begin{figure}[H]
    \centering
    \includegraphics[width=0.9\linewidth]{snow_basics/temperature_trend.png}
    \caption{Global surface temperatures anomaly relative to the 1880--1920 average. Adapted from \cite{temp_trend}.}
    \label{fig:temperature_trend}
\end{figure}
